\section{Uniform denominators and the base of a fibration}
\label{sec:denominators}

The induction needs more than a bound for the index of the general
fibre.  It also needs a bound for the Cartier denominator of the moduli
b-divisor in the canonical bundle formula.  We obtain this bound from
degenerations of the individual Beauville--Bogomolov factors.  The
absence of a horizontal boundary is essential here: it allows us to use
the form algebras of those factors and the structure-sheaf cohomology of
their degenerations.

Throughout this section, fix $n\geq1$ and assume $L_j$ for every
$j<n$: for each rational DCC coefficient set, normal projective lc
log Calabi--Yau pairs of dimension $j$ have a common principal
adjoint multiple.  We use the log abundance and good-model theorem
of the companion manuscript \cite[Theorem~11.1]{LA}.  The curve and
torsion arguments below develop the method of
\cite[Sections~3--7]{R4} in arbitrary fixed dimension.

\subsection{The exact canonical bundle formula}

\begin{proposition}[Bounded moduli denominator]\label{prop:denominator}
There are positive integers $p_0\mid p$, depending only on $n$, with
the following property.  Let $f:X\to Z$ be a contraction of normal
projective complex varieties, where $\dim X=n$, $\dim Z>0$, $X$ is
klt, and $K_X$ is rationally linearly equivalent to the pullback of
a rational Cartier divisor on $Z$.  Then there are a rational function
$\psi\in\C(X)^*$ and canonical bundle formula data satisfying the
equality of rational divisors
\begin{equation}\label{eq:den-exact}
 K_X+\frac{1}{p_0}\Div(\psi)=f^*D_Z,
 \qquad D_Z=K_Z+B_Z+M_Z.
\end{equation}
Here $B_Z\geq0$ has coefficients in a rational DCC set depending only
on $n$, the generalized pair $(Z,B_Z+M_Z)$ is generalized klt,
and the moduli b-divisor $\mathbf M$ is b-nef with $p\mathbf M$
b-Cartier.
\end{proposition}

We first construct all the data except the uniform Cartier denominator.
The geometric generic fibre is klt, has dimension less than $n$, and
has rationally trivial canonical divisor.  The induction hypothesis
therefore gives a common $p_0$ trivializing its canonical divisor.
This trivialization descends in the same degree to the generic fibre:
Cartierness descends by faithful flatness, and the one-dimensional
space of sections commutes with extension of the ground field.  A
nonzero descended section generates, since it does so after that
extension.

Choose $\psi$ so that $p_0K_X+\Div(\psi)$ has no horizontal
component.  To obtain the exact pullback in \eqref{eq:den-exact},
choose an initial rational Cartier divisor $D^0$ with
$K_X\sim_{\Q}f^*D^0$.  A sufficiently divisible multiple of
$p_0K_X+\Div(\psi)-p_0f^*D^0$ is principal and vertical.  A
rational function with vertical divisor is constant on the projective
generic fibre and hence lies in $\C(Z)$, since $f_*\OO_X=\OO_Z$.
Dividing its divisor by the chosen multiple gives the required
rational Cartier divisor $D_Z$.

Ambro's klt-trivial fibration theorem \cite{Ambro2005} now supplies
$B_Z$ and $\mathbf M$.  Its rank-one condition holds: on a
resolution, the rounded discrepancy divisor over the generic base
point is effective and exceptional, and allowing poles along an
effective exceptional divisor creates no additional rational
functions.  Thus the relevant generic direct image has rank one.
The discriminant coefficient at a prime $P\subset Z$ is
$1-t_P$, where $t_P$ is the lc threshold over its generic point of
$f^*P$.  Since the total-space boundary is zero, $0<t_P\leq1$.
Threshold ACC \cite{HMX2014} makes the resulting coefficient set a
fixed rational DCC set.  Qualitative moduli theory gives generalized
klt singularities and a smooth determination $W\to Z$ on which
$\mathbf M$ descends as a nef rational Cartier divisor.  It remains
to bound the denominators of this divisor.

\subsection{Weights of relative forms}

Fix a prime divisor $P$ on the smooth determination $W$.  Write
$\alpha$ for its coefficient in the pullback of $D_Z$, and let $t$
be the threshold of $P$ computed with the crepant total-space data.
Since canonical divisors are integral, the coefficient of $M_W$
has the same denominator as $\alpha+t$.  We calculate this number
on a transverse curve.

Here is the reduction, including its effect on actual divisors.
Resolve the total space over $W$, together with its crepant boundary
and the inverse image of $P$.  Intersect $W$ with general very ample
members until a smooth curve $C$ remains, and choose a general point
$c\in C\cap P$.  Bertini applied to the finitely many resolution
strata preserves the boundary coefficients and the SNC threshold
calculation at $c$.  The curve avoids every relevant codimension-two
exceptional subset of the base.  For each cut, use a linearly
equivalent divisor avoiding $c$ to choose the adjunction canonical
divisor: multiplication by the associated rational function followed
by residue compensates for the cut.  The compensating base divisors
have coefficient zero at $c$.  Consequently the sliced equality is
\begin{equation}\label{eq:den-curve}
 K_V+B^c+\frac{1}{p_0}\Div(\psi_C)=h^*D_C,
 \qquad \operatorname{coeff}_cD_C=\alpha,
\end{equation}
where $h:V\to C$ is the resolved slice and $B^c$ denotes its crepant
boundary.  Its geometric generic fibre is birational to a normal klt
variety with the original zero-boundary data.  This follows by making
the cuts over the open set where the resolution and the base
determination have their generic properties, and using geometric
normality and generic smoothness.  The generic fibre is geometrically
integral, so the resolved map has connected fibres.  No assertion for
a fibration with horizontal boundary is being used.

Let $z$ be a local parameter at $c$, and choose a rational top form
$\theta$ defining $K_V$.  On the generic fibre put
\[
 \phi=(\theta/dz)^{\otimes p_0}\psi_C.
\]
It generates the divisorial $p_0$-canonical sheaf of the klt generic
model.  More generally, for a relative pluricanonical form $\omega$
of degree $m$ define
\begin{equation}\label{eq:den-weight}
 w_z(\omega)=\inf_E
 \frac{\ord_E(\omega\wedge dz/z)+m}{m\ord_E(z)}.
\end{equation}
The infimum is over vertical divisorial valuations centred over $c$.
Tensor powers are understood in the wedge notation; the order in
the numerator is the ordinary canonical order on a model carrying
$E$, without a built-in logarithmic correction.

The crepant equality \eqref{eq:den-curve} gives
\begin{equation}\label{eq:den-weight-threshold}
 w_z(\phi)=\alpha-1+t.
\end{equation}
Indeed, if $b_E$ is the crepant boundary coefficient and
$v_E=\ord_E(z)$, the quotient in \eqref{eq:den-weight} is
$\alpha-1+(1-b_E)/v_E$.  Taking its infimum is precisely the
threshold definition of $t$.

\begin{lemma}[Rules for weights]\label{lem:den-weight-rules}
The weight in \eqref{eq:den-weight} has the following properties.
\begin{enumerate}
\item Finite extension of the fibre function field, with the base
field fixed, does not change the weight of a pulled-back form.
\item Under a curve extension of ramification $l$ at the marked
point, with new parameter $u$, the pulled-back weight is multiplied
by $l$.
\item Taking a positive tensor power leaves the weight unchanged.
Multiplying a degree-$m$ form by a base function $a$ adds
$\ord_z(a)/m$ to its weight.
\end{enumerate}
\end{lemma}

\begin{proof}
For a divisorial valuation of ramification index $e$ in a finite
function-field extension, the canonical ramification formula is
\[
 \ord_{E'}(\omega\wedge dz/z)+m
   =e\bigl(\ord_E(\omega\wedge dz/z)+m\bigr).
\]
With the base fixed, the denominator also multiplies by $e$.
Divisorial valuations restrict and prolong to divisorial valuations,
which proves the first assertion.  After curve ramification,
$z=u^l$ times a unit.  The logarithmic base differentials differ by
a unit, whereas $\ord_{E'}(u)=e\ord_E(z)/l$, proving the second
assertion.  One may take the curve extension Galois; regularity of
the original function field over its constant field makes the
prolongations at the different base branches conjugate, so any
chosen branch gives the same infimum.  The last assertions follow
directly from the numerator of \eqref{eq:den-weight}.
\end{proof}

\subsection{Semistable block models and their characters}

We shall use the singular Beauville--Bogomolov decomposition in the
form recalled in Lemma~\ref{lem:cover-comparison}.  After finite
extension of the curve field, the geometric generic klt fibre has a
finite quasi-\'etale product cover with one abelian block, if present,
and irreducible Calabi--Yau or irreducible symplectic blocks.  Take a
Galois comparison cover $R$ of the original fibre function field,
including all conjugates of the product cover.  Replace the curve
field by its relative algebraic closure in $\C(R)$.  This constants
extension is Galois over the original curve field, because the
original fibre function field is regular over that field.  The cover
$R$ is geometrically integral and remains geometrically quasi-\'etale
over the product.

An inertia generator on the new curve has a finite-order lift $g$ to
$R$.  By Lemma~\ref{lem:cover-comparison}, $g^b$ preserves every
block direction for $b=n!$ and induces regular semilinear
transformations on the Stein blocks $R_i$.  Each $R_i$ is a finite
quasi-\'etale cover of its original block; it therefore has the same
form-algebra properties, and has a volume generator of degree one.
Further finite Galois curve extensions cause no obstruction: take
the composite Galois function-field cover and lift the actions.
Inertia in a tower of complex curve extensions is cyclic and
surjects onto inertia below, so compatible inertia generators can
be chosen.

We record explicitly the semistable models used below.  Spread the
finite list of Stein blocks, morphisms, and actions over a punctured
curve.  Take projective equivariant resolutions; on an abelian block
do not modify its smooth generic fibre.  After one further curve
extension, projective semistable reduction \cite{KKMSD1973} gives
smooth total spaces with reduced SNC special fibres and lifted
actions of $g^b$.  Equivariance can be retained as follows.  Start
with an equivariant SNC resolution and barycentrically subdivide its
boundary complex so translates of a boundary component are equal or
disjoint.  Root extraction of the base parameter and normalization
are equivariant.  In the toroidal subdivision step, stabilized cones
then have their rays individually preserved.  Choose the compatible
height-one subdivisions on the combinatorial quotient complex and
pull them back along the orbits, using the original cone lattices,
not the lattices of a geometric quotient.  The semistable subdivision
theorem is compatible with prescribed face subdivisions; it
therefore supplies the required invariant projective subdivisions
after sufficiently divisible ramification.  This construction also
applies to strata that split after root extraction, with their
base-changed lattices.  The total spaces can be compactified and
resolved equivariantly away from the chosen special fibre.

Write $u$ for the final uniformizer, $l$ for the total ramification
over $z$, and $\zeta$ for the primitive $l$th root by which $g$
acts on the tangent line at the marked point.  On each generic Stein
block choose a volume generator $\eta_i$.  Its pullback to a smooth
resolution has no horizontal poles, by canonicality.  On a
semistable model the weight is an integer: it is the smallest
logarithmic order of $\eta_i\wedge du/u$ along a special-fibre
component.  Multiply $\eta_i$ by a power of $u$ so that
\begin{equation}\label{eq:den-block-normalization}
 w_u(\eta_i)=0.
\end{equation}
The resulting logarithmic form is regular near the fibre and has
order zero on some component.  The lc property of an SNC pair
shows that all higher divisorial tests in
\eqref{eq:den-weight} have nonnegative weight.

The product of these normalized forms also has weight zero.  To see
this without assuming that $R$ splits, first work on the fibre
product of the semistable block models.  The local equations are
$\prod_jx_{ij}=u$, one equation for each block.  Relative log top
forms multiply without an additional power of $u$.  These are
toroidal charts; their fibre products over the log base are
saturated, have reduced special fibre, and retain their log
generators on toroidal resolution.  Thus every divisorial test has
nonnegative weight.  At a product of generic points of components
on which the individual forms have order zero, all morphisms are
smooth and the product form has order zero.  Finally apply
Lemma~\ref{lem:den-weight-rules} to the finite comparison cover $R$.

The pullback of $\phi$ and the $p_0$th power of this product both
generate the generic $p_0$-canonical line on $R$. Their ratio is
therefore a base function $a$. The weight rules give
\begin{equation}\label{eq:den-product-weight}
 \phi=a\left(\bigwedge_i\eta_i\right)^{\otimes p_0},
 \qquad l\,w_z(\phi)=\frac{\ord_u(a)}{p_0}.
\end{equation}
The same formula holds for a zero-dimensional generic fibre, with
the empty product interpreted as $1$.

For each block write
\[
 (g^b)^*\eta_i=c_i\eta_i.
\]
The ratio $c_i$ is a base function, since both forms generate the
generic volume line. Their equal weights imply $\ord_u(c_i)=0$,
so $c_i$ is a unit at the marked point. Put $\lambda_i=c_i(c)$.
The finite order of the semilinear action shows, by evaluating the
product of its translates at this fixed point, that $\lambda_i$ is
a root of unity.

The form $\phi$ comes from the original function field and is
invariant under $g^b$. Taking leading coefficients in
\eqref{eq:den-product-weight} therefore gives
\begin{equation}\label{eq:den-inertia}
 \zeta^{b\ord_u(a)}\prod_i\lambda_i^{p_0}=1.
\end{equation}
Suppose a single integer $N=N(n)$ kills every $\lambda_i$. Raising
\eqref{eq:den-inertia} to the $N$th power gives
$l\mid bN\ord_u(a)$, and hence
\[
 p_0bN\,w_z(\phi)\in\Z.
\]
By \eqref{eq:den-weight-threshold}, this same integer clears the
denominator of $\alpha+t$. Thus the
remaining task is to bound the orders of the central characters
$\lambda_i$; the ramification $l$ itself need not be bounded.

\begin{lemma}[Bounded degeneration characters]\label{lem:den-characters}
In the semistable construction above, let $R_i$ be the Stein block
covers, of dimensions less than $n$, over the final curve field.
Let $c$ be the marked point, $u$ its uniformizer, and $g$ a
finite-order lift of base inertia. Put $b=n!$, so that $g^b$
preserves each block, and choose volume generators $\eta_i$ with
$w_u(\eta_i)=0$. There is an integer $N=N(n)>0$ such that, writing
\[
 (g^b)^*\eta_i=c_i\eta_i,
 \qquad \lambda_i=c_i(c),
\]
each $c_i$ is a unit at $c$ and $\lambda_i^N=1$.
\end{lemma}

\begin{proof}
\smallskip\noindent\textit{Abelian blocks.}
Let the block have dimension $s$.  On its projective semistable
model over a disc, the direct image of the relative log top line is
the extended Hodge line $F^s$ in degree $s$.  This is the semistable
comparison theorem for logarithmic de Rham cohomology and the
canonical extension with nilpotent residue
\cite{Deligne1970,Steenbrink1976}.  The normalized form is a local
frame: it is a regular relative log form but cannot be divided by
$u$ while remaining regular, because its order on one component is
zero.  Its eigenvalue on the central fibre of $F^s$ is $\lambda_i$.

The integral local system in degree $s$ has rank
$\binom{2s}{s}$.  Choose a coordinate in which the finite base
action is a rotation.  On the universal cover of the punctured disc,
the lifted transformation, combined with parallel transport, acts
by an integral matrix $A$ commuting with the unipotent monodromy.
A power of $A$ is a power of monodromy, because the original
transformation has finite order.  Passing from flat frames to
canonical-extension frames multiplies $A$ by the exponential of a
scalar multiple of the logarithm of monodromy.  This is a commuting
unipotent factor, so it does not change the eigenvalues.  Hence a
primitive root of the order of $\lambda_i$ occurs among the
eigenvalues of an integral matrix of rank $\binom{2s}{s}$.  Its
cyclotomic degree is at most that rank.  There are only finitely many
possible orders, uniformly for $s<n$.

\smallskip\noindent\textit{Nonabelian blocks: a good special-fibre model.}
Let $s$ be the block dimension, write $\eta$ for its normalized
form, and let $G$ be the finite cyclic group generated by the
transformation being considered. We will bound the character by
applying Lemma~\ref{lem:finite-character} to a normal component
of the special fibre. This requires a degree-one log generator on
that component and a bounded power of the transformation preserving
it. We first construct an equivariant good model on which the
normalized logarithmic form generates along the whole special
fibre; we then use structure-sheaf cohomology to find the required
power. We claim that, near the marked
point, there is a $G$-equivariant projective model $Y$ such that
\begin{equation}\label{eq:den-good-model}
 Y\text{ is klt},\qquad T=Y_c\text{ is reduced Cartier},
 \qquad (Y,T)\text{ is dlt},\qquad K_Y\text{ is semiample over the curve}.
\end{equation}
The geometric generic fibre need only remain birational to the
block.

Start from an equivariant smooth projective compactification $Q$ of
the semistable model, preserving its SNC special fibre.  A relative
$K_Q$-MMP is also a $(K_Q+Q_c)$-MMP, because $Q_c$ is a fibre.
For existence and termination, take the finite quotient
$\overline Q=Q/G$ and its branch boundary $\Delta$, so that
$K_Q$ is the finite pullback of $K_{\overline Q}+\Delta$.
The quotient is $\Q$-factorial and $(\overline Q,\Delta)$ is
klt.  Its adjoint is pseudo-effective over the quotient curve:
an effective pluricanonical form on the generic fibre upstairs
gives an invariant one after multiplying its finitely many
translates.

Choose a divisor $H$ of sufficiently large degree on the quotient
curve, represented by many general points with small rational
coefficients.  The pair obtained by adding its pullback stays klt,
and its adjoint is pseudo-effective absolutely, since a generic
section extends after a sufficiently positive base twist.  Choose
also $\deg H>2\dim Q$.  The good-model theorem of
\cite[Theorem~11.1]{LA}, followed by a terminating MMP with scaling
to that good model, now applies.  Every step is over the curve:
otherwise the extremal-ray length bound for
$(\overline Q,\Delta)$ would give a horizontal generating curve
$C_0$ with
$-(K_{\overline Q}+\Delta)\cdot C_0\leq2\dim Q$, while
the pullback of $H$ has degree greater than this on $C_0$.
Inductively the base pullback and the morphism to the curve persist.

Lift each step by normalization in the upstairs function field,
using Stein factorization for the contractions.  Finite pullback
preserves the relative negativity and ampleness needed for flips.
No divisors are extracted, and the branch-canonical equality remains
valid in codimension one.  This gives the finite-group equivariant
MMP; invariant $\Q$-factoriality suffices upstairs
\cite[Theorems~3.3.2 and~3.4.2]{Prokhorov2021}.  The steps factor
over the upstairs curve as well, by normalization in its function
field.  They preserve dlt for the pair with the special fibre added.
Special-fibre multiplicities remain one because there is no
extraction.  The klt total space is Cohen--Macaulay, so its Cartier
fibre has no embedded associated components; generic reducedness
therefore makes the fibre reduced scheme-theoretically.
Relative semiampleness pulls back from the
downstairs good model.  This proves \eqref{eq:den-good-model}.

\smallskip\noindent\textit{The degree-one log generator.}
The logarithmic divisor of $\eta\wedge du/u$ on $Y$ is effective
near $T$: this held on the semistable model, and the program extracts
no divisors.  The generic fibre has Kodaira dimension zero.
Consequently the semiample contraction of $K_Y$ has
zero-dimensional image on the generic fibre; its Stein image over
the curve is the curve itself.  Thus $K_Y$, and also $K_Y+T$, is
rationally linearly trivial over a neighbourhood of $c$.
The effective log divisor has no horizontal part, since its
restriction to the generic fibre is effective and rationally
trivial.  A principalization of a multiple uses a rational function
with vertical divisor, hence a base function.  The log divisor is
therefore $aT$ for some $a\in\Q_{\geq0}$.  If $a>0$, every
weight quotient is at least $a$: it is $a$ plus the log discrepancy
of $(Y,T)$ divided by the fibre multiplicity.  This contradicts
$w_u(\eta)=0$.  Hence
\begin{equation}\label{eq:den-log-generator}
 \Div_{\log}(\eta\wedge du/u)=0.
\end{equation}
In particular $K_Y+T$, and therefore $K_Y$, is Cartier near $T$.

Adjunction makes $T$ Gorenstein with trivial dualizing line, generated
by the residue of \eqref{eq:den-log-generator}.  The transformation
acts on that generator by $\lambda$: the residue of its action on
$du/u$ is one.  If a prime component $S\subset T$ is preserved by
a power of the transformation, the residue on $S$ generates
$K_S+\Theta_S$, where $\Theta_S$ is its different.  Dlt
adjunction makes $(S,\Theta_S)$ normal lc and effective.  Moreover
$\Theta_S$ is integral, since the divisor of this degree-one
rational form is integral and equals $-\Theta_S$.  Its coefficients
are at most one, so it is reduced.  The character on this component
generator is the corresponding power of $\lambda$.

Thus a bounded power preserving one component will reduce the
character bound to Lemma~\ref{lem:finite-character}. We obtain
such a power from the structure-sheaf cohomology of $T$.

\smallskip\noindent\textit{Finding a preserved component.}
The family is flat over the smooth curve.  Its central fibre is slc
by \eqref{eq:den-good-model}, and its nearby fibres are klt after
shrinking; these fibres are Du Bois.  Cohomology and base change for
proper flat Du Bois families makes $h^j(\OO_{Y_t})$ locally
constant \cite[Corollary~2.64]{KollarFamilies2023}.  The nearby fibres are
birational to the block and have rational singularities.  Their
structure-sheaf cohomology is therefore that of a smooth resolution
of the block.  Extension of differential forms on klt spaces,
together with Hodge symmetry, identifies these dimensions with the
form algebra of the block \cite{GKKP2011}.

For an irreducible Calabi--Yau block, the only nonzero groups
$H^j(T,\OO_T)$ occur in degrees $0$ and $s$, and each has
dimension one.  Serre duality and the log generator give alternating
trace
\[
 1+(-1)^s\lambda^{-1}.
\]
If $\lambda$ has order greater than two, this is nonzero.  In the
symplectic case there is one dimension in each even degree and none
in odd degree.  For some power of the transformation of exponent at
most $s+1$ the trace is nonzero.  Indeed, if the power sums of the
finitely many nonzero eigenvalues all vanished through their number,
Newton's identities would force their product to vanish.

The coherent Lefschetz theorem implies that a finite-order
transformation with nonzero alternating trace has a fixed point
\cite{Donovan1969}.  Its vanishing statement when there are no fixed
points applies to the possibly singular projective scheme $T$, for
example by equivariant embedding and coherent Lefschetz localization.
At a fixed point at most $s+1$ components of $T$ meet: intersections
of components of a dlt boundary are lc strata, with the generic SNC
codimension description.  A further power of exponent dividing
$(s+1)!$ therefore preserves a prime component through that point.
Its reduced log Calabi--Yau pair has dimension $s<n$ and the
degree-one generator just constructed.  Lemma~\ref{lem:finite-character},
which follows from $L_s$, bounds the order of its character.
This bounds the order of $\lambda$ as well.  Taking a common
multiple over all $s<n$ and both types of block proves the lemma.
\end{proof}

\subsection{Completion of the denominator bound}

Choose $N$ from Lemma~\ref{lem:den-characters} and set $p=p_0bN$.
The calculation preceding the lemma shows that $p$ clears
$\alpha+t$ at every prime divisor of $W$. Thus every coefficient
of $pM_W$ is integral. Since $W$ is smooth, this divisor is
Cartier. The b-divisor already descends to $W$ by qualitative
moduli theory, so $p\mathbf M$ is b-Cartier. This completes the
proof of Proposition~\ref{prop:denominator}.

\subsection{Big bases and their complete section fields}

\begin{corollary}\label{cor:big-base}
Under the hypotheses of Proposition~\ref{prop:denominator}, suppose
that $D_Z$ is big.  There is a positive integer $m=m(n)$, divisible
by $p_0$, such that $|mK_X|$ is nonempty and its section-ratio field
inside $\C(X)$ is exactly $\C(Z)$.
\end{corollary}

\begin{proof}
Take a smooth determination $\tau:W\to Z$ on which $pM_W$ is
nef Cartier, and resolve the boundary support.  Let $\Gamma$ be
the strict transform of $B_Z$ plus the reduced exceptional divisor.
Generalized klt gives
\[
 K_W+\Gamma+M_W=\tau^*D_Z+E,
 \qquad E\geq0\text{ exceptional}.
\]
This is a big adjoint for a log smooth lc pair whose coefficients
belong to a fixed DCC set, with fixed Cartier denominator for the
nef part.  Polarized effective birationality
\cite[Theorem~1.3]{BirkarZhang2016} gives a uniform degree whose
rounded system is birational.  Such sections push forward as
rational sections of the rounded system of $D_Z$, because the
two divisors agree at the nonexceptional generic divisorial points.
Thus a uniform rounded system downstairs is birational.  Replacing
its degree by a common multiple preserves its section ratios and
makes the degree divisible by $p_0$.

For every $m$ divisible by $p_0$, the exact identity is
\begin{equation}\label{eq:den-section-comparison}
 H^0(X,\OO_X(mK_X))
 =\psi^{m/p_0}f^*
 H^0\bigl(Z,\OO_Z(\lfloor mD_Z\rfloor)\bigr),
\end{equation}
where both sides are viewed as divisorial spaces of rational
functions with compatible canonical representatives.  One
inclusion follows by pulling back an effective rational Cartier
divisor.  For the converse, divide an upstairs section by the common
factor $\psi^{m/p_0}$.  Its restriction to the projective generic
fibre is a regular function, hence belongs to $\C(Z)$.  Checking
orders over each prime of $Z$ shows that
$\Div(h)+mD_Z\geq0$, equivalently
$\Div(h)+\lfloor mD_Z\rfloor\geq0$.  This proves
\eqref{eq:den-section-comparison}.  The birational system on $Z$
then gives precisely its full function field, not merely a finite
subfield.
\end{proof}

\subsection{Torsion over rationally connected bases}

The second application requires an integral principal multiple of
the actual generalized adjoint, not merely a bound for its Cartier
index.  We give the general statement, since it also handles the
ordinary rationally connected log Calabi--Yau case.

\begin{proposition}\label{prop:rc-torsion}
Fix a dimension $d$, a rational DCC set $I\subset[0,1]$, and a
positive integer $p$.  There is an integer $\ell=\ell(d,I,p)>0$
with the following property.  Let $Z$ be a normal projective complex
variety of dimension $d$ with a rationally connected smooth
projective resolution.  Suppose $(Z,B+M_Z)$ is generalized klt,
$B\geq0$ has coefficients in $I$, the b-nef data satisfy
$p\mathbf M$ b-Cartier, and
\[
 D=K_Z+B+M_Z\sim_{\Q}0.
\]
Then $\ell D$ is an integral principal divisor.
\end{proposition}

\begin{proof}
We follow the generalized torsion argument of
\cite[Proposition~7.1]{R4}, keeping track of why no
dimension-four restriction is needed.

We need two uniform integers: one making $D$ integral, and one
killing the resulting torsion class in the Weil divisor class
group $\Cl(Z)$. Global ACC, after the extractions below, gives a
finite coefficient set and a uniform positive lower bound for
generalized log discrepancies. The latter permits boundedness;
topology of the regular loci and Kummer theory then give a common
torsion exponent.

\smallskip\noindent\textit{Extraction and global ACC.}
We first justify the small modifications and one-divisor extractions
used with generalized global ACC.  Choose a high log resolution
$g:W\to Z$ on which the nef data descend and, if a valuation $E$
is marked, on which $E$ appears.  Write
\[
 K_W+B_W+M_W=g^*D.
\]
The boundary $B_W$ is sub-klt and need not be effective.  There is
an effective exceptional Cartier divisor $F$ with $-F$ relatively
ample; construct the resolution as a sequence of blowups and take
suitable positive combinations of the successive exceptional
divisors.  Resolve all supports.  If $E$ has generalized log
discrepancy $a\in(0,1)$, choose $\delta>0$ small enough that
$B_W+\delta F$ remains sub-klt and
$\delta\operatorname{mult}_EF<a$.  For a sufficiently ample
Cartier divisor $A$ on $Z$, choose a general effective rational
divisor
\[
 \Lambda\sim_{\Q}M_W+g^*A-\delta F
\]
by dividing a general very ample member.  Then
$(W,B_W+\delta F+\Lambda)$ is sub-klt.  Pushing down and
comparing the rational principal divisors shows that
$\Xi=B+g_*\Lambda$ is an effective ordinary klt boundary on
$Z$.  Its discrepancy at the marked $E$ is
$a-\delta\operatorname{mult}_EF\in(0,1)$.

The extraction theorem \cite[Corollary~1.4.3]{BCHM2010} applied to
$(Z,\Xi)$ extracts exactly $E$ when $E$ is exceptional, or gives a
small $\Q$-factorialization if the marked list is empty.  If the
marked valuation is already a prime on $Z$, use the empty list.
Restore the
original generalized data on this model $Z^+$.  The boundary is
effective: it is the strict transform of $B$, together with
$(1-a)E$ in the extraction case.  The generalized adjoint remains
rationally trivial and the nef denominator remains $p$.

If the nef data are numerically nontrivial, write them as the single
nef Cartier term $pM_W$ with weight $1/p$.  Generalized global ACC
\cite[Theorem~1.6]{BirkarZhang2016} applies.  If they are
numerically trivial, take a common higher model
$\pi:W'\to Z^+$ determining the nef data.  Since $Z^+$ is
$\Q$-factorial, its nef trace is rational Cartier.  The difference
$M_{W'}-\pi^*M_{Z^+}$ is exceptional and numerically trivial over
$Z^+$.  Negativity applied to it and its negative shows that the
difference is zero.  The generalized discrepancies are then
ordinary discrepancies, and ordinary global ACC
\cite[Theorem~1.5]{HMX2014} applies instead.  Apply this discussion
first with no marked valuation.  It places the original boundary
coefficients in a fixed finite subset $I_0$ of $I$.

The same argument gives a uniform $\epsilon>0$ for generalized
$\epsilon$-lc singularities.  Otherwise choose marked valuations
whose discrepancies decrease strictly to zero.  Extract them as
above.  The new coefficients $1-a_i$ form a strictly increasing
sequence, so adjoining them to $I$ still gives a DCC set.  Global
ACC would force them into a finite set, a contradiction.

\smallskip\noindent\textit{Bounded topology and class-group torsion.}
Birkar's boundedness theorem \cite[Theorem~1.7]{Birkar2025} applies
in every fixed dimension: the varieties are projective and
rationally connected, their generalized adjoints are real-linearly
trivial, and they are uniformly generalized $\epsilon$-lc.
Consequently they are bounded up to isomorphism in codimension one.
The bounded comparison models can be normalized, still in a bounded
family and still isomorphic to $Z$ in codimension one.

Put $U=Z_{\mathrm{reg}}$.  The groups $H_1(U(\C),\Z)$ are
finitely generated and range over only finitely many isomorphism
types.  First, this group is invariant under isomorphism in
codimension one of normal varieties.  On the regular loci the
common open is obtained by removing subsets of complex codimension
at least two.  After stratifying those subsets, transversality moves
paths and their homotopies off every stratum, whose real codimension
is at least four; the fundamental group is unchanged.  Second,
stratify a finite-type parameter space for the bounded comparison
family until it is flat.  The relative smooth locus then has fibres
equal to the regular loci of the selected normal fibres.  Cover the
parameter strata by finitely many affine opens and embed the family
projectively over each.  Complex projective space is a compact real
algebraic set via its realization as rank-one Hermitian projection
matrices.  Thus the entire relative smooth locus, not just individual
affine fibre charts, maps semialgebraically to each parameter open.
Semialgebraic triviality gives a finite partition on which these
fibres are homeomorphic, and semialgebraic triangulation gives each
the homotopy type of a finite complex \cite{DelfsKnebusch1982}.
This proves the asserted finite list of finitely generated homology
groups.

For each individual $Z$, the class group $\Cl(Z)$ is finitely
generated.  Indeed a rationally connected smooth resolution $Y$
has $\operatorname{Pic}^0(Y)=0$, so $\operatorname{Pic}(Y)$ is
finitely generated by Neron--Severi finiteness; divisor pushforward
surjects onto $\Cl(Z)$.  Since $Z$ is normal and projective,
$\Gamma(U,\OO_U^*)=\C^*$ and
$\operatorname{Pic}(U)=\Cl(Z)$.  Kummer theory and comparison
for finite etale covers \cite{SGA1} give, for every $a\geq1$,
\begin{equation}\label{eq:den-kummer}
 \Cl(Z)[a]\simeq H^1_{\mathrm{\acute et}}(U,\mu_a)
 \simeq\operatorname{Hom}\bigl(H_1(U(\C),\Z),\mu_a(\C)\bigr).
\end{equation}
The left side has bounded order as $a$ varies for this fixed $Z$.
Thus $H_1(U(\C),\Z)$ has no free summand and is finite.  There
are only finitely many such finite groups in our bounded topological
list.  A common multiple $T$ of their exponents kills every torsion
class in every $\Cl(Z)$, by \eqref{eq:den-kummer}.

\smallskip\noindent\textit{Clearing the actual divisor.}
The trace $pM_Z$ is an integral Weil divisor, being the pushforward
of Cartier data.  Choose $q$ divisible by $p$ and by the
denominators of the finite set $I_0$.  Then $qD$ is integral.
Its class in $\Cl(Z)$ is torsion because $D\sim_{\Q}0$.
Therefore $TqD$ is principal.  Taking $\ell=Tq$ proves the
proposition.
\end{proof}

\begin{corollary}\label{cor:relative-rc-torsion}
Under the hypotheses of Proposition~\ref{prop:denominator}, suppose
$K_X\sim_{\Q}0$ and a smooth projective resolution of $Z$ is
rationally connected.  There is an integer $r=r(n)>0$ such that
$rK_X$ is an integral principal divisor.
\end{corollary}

\begin{proof}
The divisor $D_Z$ in \eqref{eq:den-exact} is rationally linearly
trivial.  Indeed a Cartier multiple whose pullback is trivial is
itself trivial by the projection formula and
$f_*\OO_X=\OO_Z$.  Proposition~\ref{prop:rc-torsion} gives a
uniform $\ell$ and a function $v\in\C(Z)^*$ with
$\ell D_Z=\Div(v)$.  Choose a common multiple $r$ of $\ell$
and $p_0$.  Then
\[
 rK_X=\Div\left((v\circ f)^{r/\ell}\psi^{-r/p_0}\right).
\]
All constants depend only on $n$, since the dimensions of $Z$
range over a fixed finite set and Proposition~\ref{prop:denominator}
supplies a fixed coefficient set and nef denominator.
\end{proof}
